\chapter{2020 Nobel Prize in Economics}
\textbf{Laureates: }Paul Milgrom (USA), Robert Wilson (USA)

\section*{Reason for Award}
For improvements to auction theory and inventions of new auction formats

\section{What They Did}
Paul Milgrom and Robert Wilson advanced auction theory and designed new auction formats that
have been used to raise billions of dollars for governments worldwide. Their work combined
theoretical insights with practical design to create auctions that are efficient, revenue-
maximizing, and resistant to manipulation.

Wilson developed the theory of common value auctions, where the item being sold has the same
value to all bidders but bidders receive different signals about that value. He analyzed the
winner's curse--the tendency for the winner of a common value auction to overpay because they
had the most optimistic signal. Wilson showed how bidders should shade their bids to account
for the winner's curse and developed the first systematic theory of how bidders behave in
common value settings. His work on the theory of bidding rings analyzed how cartels can
undermine auction efficiency.

Milgrom and Wilson developed the theory of combinatorial auctions, where bidders can bid on
packages of items rather than individual items. This is important when items are complements
and the value of a package exceeds the sum of individual values. Their work showed how to
design auctions that efficiently allocate complementary items while preventing collusion.

They designed the simultaneous multiple round auction (SMRA) used by the Federal Communications
Commission (FCC) to auction radio spectrum licenses. This auction format has raised tens of
billions of dollars for the US Treasury. Milgrom and Wilson also designed the combinatorial
clock auction used for spectrum auctions in several countries. The clock auction allows
bidders to express preferences over packages of items while maintaining transparency and
efficiency.

Milgrom's earlier work on auction theory included the revenue equivalence theorem and the
theory of auction bidding strategies. He showed how different auction formats affect bidder
behavior and seller revenue. His work on dynamic auctions analyzed how sellers should set
reserve prices and manage bidding processes over time.

Wilson also contributed to the theory of bidding behavior in auctions with affiliated
values. He showed how bidders' signals are correlated and how this affects equilibrium
bidding strategies. His work on the theory of sequential auctions analyzed how the order
of sales affects outcomes.

\section{Impact on Economic Thinking}
Their work transformed auction theory from an abstract mathematical discipline into a practical
tool for designing efficient markets. The theoretical insights they developed have been applied
to a wide range of auction contexts, from spectrum auctions to electricity markets to online
advertising.

The practical applications of their work have been enormous. Spectrum auctions designed using
their insights have raised hundreds of billions of dollars worldwide, funding public services
and infrastructure. Their work has also influenced the design of auctions for other goods
and services, from emissions permits to procurement contracts.

The theoretical framework they developed has become standard in auction theory and mechanism
design. Their insights about the winner's curse, complementary values, and auction design
have been widely applied in both academic research and practical auction design. The common
value model has been extended to analyze bidding in oil and gas leases, mineral rights, and
other natural resource auctions.

Wilson's analysis of the winner's curse has influenced how companies approach bidding in
acquisition contexts and other common value settings. Milgrom's work on dynamic auctions
has been applied to online advertising platforms like Google and Facebook, which run billions
of dollars of auctions daily.

Their work demonstrated that auction design matters for economic outcomes. By carefully
structuring auction rules, designers can achieve better outcomes for both sellers and
buyers. This insight has been applied to diverse settings including spectrum allocation,
airline slot allocation, and kidney exchange programs.

\section{Modern Perspectives Today}
Auction theory based on Milgrom and Wilson's work continues to evolve, with new applications
to online advertising, electricity markets, and other contexts. The combinatorial clock auction
has been used for spectrum auctions in Europe and Asia.

The theoretical framework continues to be extended to new contexts, including dynamic auctions,
network industries, and platform markets. Their work remains foundational for understanding
how to design efficient auctions in a wide range of contexts. Research on algorithmic auction
design has been particularly important for digital platforms.

The practical applications of their work continue to expand, with new auction formats being
designed for emerging markets like carbon emissions trading and bandwidth allocation. Their
contributions to auction theory have had a profound impact on both economic theory and practical
market design. The FCC auctions designed by Milgrom and Wilson have become models for spectrum
auctions worldwide.

Recent work has extended their framework to analyze auctions with budget constraints, risk-
averse bidders, and heterogeneous values. These extensions have been important for designing
auctions in developing countries and for emerging technologies like 5G spectrum and satellite
internet. Their theoretical insights continue to inform auction design for internet advertising
and other digital markets.
